\documentclass[10pt,a4paper]{amsart}
\usepackage[margin=19mm]{geometry}
\usepackage{amsmath,amssymb,mathtools}
\usepackage[scr=rsfs,cal=euler]{mathalfa}
\usepackage{xfrac}
\usepackage[unicode,hidelinks,bookmarks=false]{hyperref}
\newcommand{\ie}{i.e.\ }
\newcommand{\cf}{cf.\ }
\newcommand{\resp}{respectively\ }
\newcommand{\Subsection}{\subsection}
\providecommand{\nobreakdash}{\nobreak\hbox{-}}
\setlength{\emergencystretch}{2em}
\multlinegap0pt
\usepackage{amssymb}
\usepackage{mathtools}
\usepackage[shortlabels]{enumitem}
\newtheorem{lem}{Lemma}
\newcommand{\M}{\operatorname{M}}
\title{Quantitative Oppenheim in signature $(2,2)$ \\ via determinant values}
\author{Wooyeon Kim, Hee Oh}
\date{Source: arXiv:2607.18037v2 (CC BY 4.0)}
\begin{document}
\maketitle

\noindent\emph{Source and licence.} Quantitative Oppenheim in signature $(2,2)$ \\ via determinant values. Version arXiv:2607.18037v2 (CC BY 4.0), distributed by arXiv under the Creative Commons Attribution 4.0 International licence, http://creativecommons.org/licenses/by/4.0/, as recorded in the arXiv licence metadata for this exact version and shown on the abstract page https://arxiv.org/abs/2607.18037v2. Full licence notice: \texttt{licenses/arxiv-2607.18037-CC-BY-4.0.txt}.\par\smallskip

\noindent\textit{Editorial scope.} Counted unit: the article's lem lem:d2-reconstruction-from-quasinull (source0.tex lines 809--822). The article's complete original proof of it is delivered with it (source0.tex lines 825--1040, one \texttt{proof} environment of 216 lines). No mathematics is written, repaired or re-derived: the statement and the proof are the article's own lines, in the article's own notation, and no retained line is substituted or re-flowed.  No further source-local context is needed: every cross-reference of every retained line resolves inside the two delivered spans.  Nothing was omitted from any retained span, so the package carries no omission record.  No retained line contains a citation, so the wrapper carries no bibliography and no bibliography entry.  No retained line contains a figure environment, an image-inclusion command or drawing code, so no image is delivered and the package declares no asset dependency.  Editorial formatting only, in addition: an AMS \texttt{amsart} preamble that restates the article's own declarations and macros used by the retained lines, namely the environments \texttt{lem} on separate counters, together with the standard packages those lines need.  The retained environments print in this document as Lemma 1 in order of appearance, whereas the article prints the same items with the numbering of its own full document.  The complete native source of the article as distributed in the arXiv e-print for this exact version is bundled unchanged as \texttt{sources/205621/source0.tex}.
\begin{lem}[Reconstruction from five nearly isotropic planes]
\label{lem:d2-reconstruction-from-quasinull}
Let $\Lambda<\M_2(\mathbb R)$ be a lattice and let $\mathcal B$ be a
$\mathbb Z$-basis of $\Lambda$. There exists $C=C(\mathcal B)>0$ such
that, for every $R\ge1$ and $0<\varepsilon<1$, if
$\#\mathcal Q_\Lambda(R,\varepsilon)\ge5$, then there exists a rational
split quadratic form $P$ satisfying
\begin{equation}\label{eq:rational-split-reconstruction}
\operatorname{ht}(P)\le CR^{18},
\qquad
\operatorname{dist}([Q_{\Lambda,\mathcal B}],[P])
\le CR^{18}\varepsilon.
\end{equation}
\end{lem}

\begin{proof}
 After increasing $C$, we
may assume that
$
\varepsilon\le cR^{-18}
$
for a sufficiently small constant $c>0$; otherwise the conclusion follows
by taking any fixed rational split form.

Choose five distinct planes in $\mathcal Q_\Lambda(R,\varepsilon)$. By the
pigeonhole principle, three of them, denoted by $V_1,V_2,V_3$, have
primitive Pl\"ucker vectors $w_1,w_2,w_3$ satisfying
$$
\|w_i-\pi_cw_i\|\le\varepsilon
\qquad
(1\le i\le3),
$$
or the analogous inequalities with $\pi_r$. We treat the first case; the
second is identical.

Put $\widehat w_i=w_i/\|w_i\|$. Since $\wedge^2\Lambda$ is a fixed lattice,
$\|w_i\|\gg1$. We claim that there is a unit Pl\"ucker vector
$z_i\in\mathcal M_c$ such that
\begin{equation}\label{eq:near-exact-ruling}
\|\widehat w_i-z_i\|\ll\varepsilon.
\end{equation}
To see this, let
$
e_1=E_{11}, e_2=E_{12},
e_3=E_{21}, e_4=E_{22},
e_{ij}=e_i\wedge e_j,
$
and write
$$
\widehat w_i
=
\alpha e_{12}
+\beta(e_{14}-e_{23})
+\gamma e_{34}
+\rho e_{13}
+\sigma(e_{14}+e_{23})
+\tau e_{24}.
$$
The Pl\"ucker relation gives
$$
\alpha\gamma-\beta^2=\rho\tau-\sigma^2.
$$
The last three coefficients are $O(\varepsilon)$, and hence
$|\alpha\gamma-\beta^2|\ll\varepsilon^2$. Moreover, one of
$|\alpha|,|\beta|,|\gamma|$ is bounded below. If $|\alpha|\gg1$, replace
$\gamma$ by $\beta^2/\alpha$; if $|\gamma|\gg1$, replace $\alpha$ by
$\beta^2/\gamma$. In the remaining case, $|\beta|\gg1$ and
$\alpha\gamma>0$, so we replace $\beta$ by
$\operatorname{sgn}(\beta)\sqrt{\alpha\gamma}$. In each case the change is
$O(\varepsilon^2)$, and the resulting nonzero vector in $\mathcal M_c$ is
a Pl\"ucker vector. Normalizing it proves
\eqref{eq:near-exact-ruling}.

We next show that the planes $V_i$ are pairwise transverse. In the integral
coordinates induced by $\mathcal B$, the vectors $w_i$ are primitive
integral vectors of norm $O(R)$. Since $V_i\ne V_j$, the vectors $w_i$ and
$w_j$ are not proportional. Some $2\times2$ minor formed from their
coordinates is therefore a nonzero integer, and hence
\begin{equation}\label{eq:plucker-separation}
\min_{\pm}\|\widehat w_i\pm\widehat w_j\|
\gg R^{-2}
\qquad
(i\ne j).
\end{equation}
By \eqref{eq:near-exact-ruling}, the same lower bound, up to a fixed
constant, holds with $z_i,z_j$ in place of
$\widehat w_i,\widehat w_j$.

Every unit Pl\"ucker vector in $\mathcal M_c$ has, up to sign, the form
$$
z(s,t)
=
(se_1+te_3)\wedge(se_2+te_4),
\qquad
s^2+t^2=1.
$$
A direct calculation gives
$$
|z(s,t)\wedge z(s',t')|
=
(st'-ts')^2
=
\frac12\min_{\pm}\|z(s,t)\pm z(s',t')\|^2.
$$
It follows that $|z_i\wedge z_j|\gg R^{-4}$. Since
$\varepsilon\ll R^{-18}$, equation \eqref{eq:near-exact-ruling} now implies
that $w_i\wedge w_j\ne0$. Thus
\begin{equation}\label{eq:pairwise-transverse}
V_i\cap V_j=\{0\}
\qquad
(i\ne j).
\end{equation}

Let $g_{\mathcal B}:\mathbb R^4\to\M_2(\mathbb R)$ send the standard basis
to $\mathcal B$. Pull back the planes $V_i$ and the planes represented by
$z_i$ under $g_{\mathcal B}$, and retain the same notation.  In these
coordinates, $\Lambda=\mathbb Z^4$ and
$$
Q:=\det\circ g_{\mathcal B}=Q_{\Lambda,\mathcal B}.
$$
By Minkowski's second theorem, for each $i$ there are linearly independent
vectors $u_i,v_i\in V_i\cap\mathbb Z^4$ satisfying
$$
\|u_i\|,\|v_i\|\ll R.
$$
Let $L$ be the $9\times10$ integer matrix representing the nine linear
functionals
$$
P\mapsto P(u_i),\quad
P(v_i),\quad
P(u_i+v_i)-P(u_i)-P(v_i),
\qquad
1\le i\le3,
$$
on the ten-dimensional space $\operatorname{Sym}^2((\mathbb R^4)^*)$ of quadratic forms on $\mathbb R^4$. Then
$$
\ker L
=
\{P:P|_{V_1}=P|_{V_2}=P|_{V_3}=0\},
$$
and every entry of $L$ is $O(R^2)$. We claim that $\ker L$ is one-dimensional. By
\eqref{eq:pairwise-transverse}, we have
$\mathbb R^4=V_1\oplus V_2$, and $V_3$ is the graph of an invertible linear
map $T:V_1\to V_2$. A quadratic form vanishing on both $V_1$ and $V_2$ has
the form
$$
P(x+y)=2\beta(x,y)
\qquad
(x\in V_1,\ y\in V_2)
$$
for some bilinear form $\beta:V_1\times V_2\to\mathbb R$. It also vanishes
on $V_3$ precisely when $\beta(x,Tx)=0$ for every $x\in V_1$, or
equivalently when the form
$$
(x,y)\mapsto\beta(x,Ty)
$$
is alternating. Since the space of alternating forms on the
two-dimensional space $V_1$ is one-dimensional, $\dim\ker L=1$.
The rational line $\ker L$ is generated by a nondegenerate rational
split form. Let $Q_0$ be its primitive integral generator.
 Its
coefficients are, up to a common divisor, the signed $9\times9$ minors of
$L$. Since the entries of $L$ are $O(R^2)$,
\begin{equation}\label{eq:Q0-height}
\operatorname{ht}(Q_0)\ll R^{18}.
\end{equation}

It remains to compare $Q_0$ with $Q$. Let $W_i$ be the pulled-back
isotropic plane represented by $z_i$. The distance between the orthogonal
projections onto two planes is controlled by the distance between their
normalized Pl\"ucker vectors. Hence, if $P_U$ denotes orthogonal projection
onto $U$,
$$
\|P_{V_i}-P_{W_i}\|_{\operatorname{op}}
\ll
\|\widehat w_i-z_i\|
\ll\varepsilon.
$$
Since $Q$ vanishes on $W_i$, it follows that
$
\left|\frac{Q(x)}{\|Q\|}\right|
\ll
\varepsilon\|x\|^2$ for $x\in V_i$.
Applying this estimate to $u_i$, $v_i$, and $u_i+v_i$ gives
\begin{equation}\label{eq:LQ-small}
\left\|
L\left(\tfrac{Q}{\|Q\|}\right)
\right\|
\ll
R^2\varepsilon.
\end{equation}

Let $\sigma_1\ge\cdots\ge\sigma_9>0$ be the nonzero singular values of
$L$. For each set $I$ of nine columns, let $L_I$ be the corresponding
$9\times9$ submatrix. The Cauchy--Binet formula gives
$$
\det(LL^{\mathsf T})
=
\sum_{|I|=9}\det(L_I)^2.
$$
Since $L$ has rank $9$, one of these minors is a nonzero integer. Therefore
$$
1
\le
\det(LL^{\mathsf T})
=
\prod_{j=1}^9\sigma_j^2.
$$
As $\sigma_j\ll R^2$, we obtain $\sigma_9\gg R^{-16}$. It follows from
\eqref{eq:LQ-small} that
$$
\operatorname{dist}
\left(
\tfrac{Q}{\|Q\|},\ker L
\right)
\le
\sigma_9^{-1}
\left\|
L\left(\tfrac{Q}{\|Q\|}\right)
\right\|
\ll
R^{18}\varepsilon.
$$
Since $\ker L=\mathbb RQ_0$, this gives
$
\operatorname{dist}([Q],[Q_0])
\ll
R^{18}\varepsilon.
$
Together with \eqref{eq:Q0-height}, this proves the lemma.
\end{proof}

\end{document}
